% Source: https://artofproblemsolving.com/community/c3h3163958_2013_emcc_individual_accuracy_test__exeter_math_club_competition
% Stable identity: aops-topic-3163958-q10
Put $u=a+b,v=b+c,w=c+a$. The inverse formulas are $a=\dfrac{u+w-v}{2}$, $b=\dfrac{u+v-w}{2}$, $c=\dfrac{v+w-u}{2}$, so integrality is equivalent to $u+v+w$ even. Since $uvw=1800$ is even, at least one factor is even. The additional requirement that $u+v+w$ be even permits exactly one or three even factors; both cases must be counted. For positive absolute values, distribute the exponents of $1800=2^3 3^2 5^2$ among three factors. An even sum requires either all three factors even, giving one allocation of the three 2 s, or just one factor even, giving three allocations. Each odd-prime exponent has $\binom42=6$ allocations. Thus there are $(1+3)6^2=144$ positive triples of factors. Each has four sign patterns with positive product, and parity is unchanged by signs. The inverse map is bijective, giving $\boxed{576}$.
